\chapter*{ТЕРМИНЫ И ОПРЕДЕЛЕНИЯ}
\addcontentsline{toc}{chapter}{ТЕРМИНЫ И ОПРЕДЕЛЕНИЯ}
\setlength{\parindent}{0pt}
% Перечень располагается столбцом, слева без абзацного отступа в алфавитном порядке,
% справа через тире — определение, без знаков препинания в конце строки (п. 4.9.2)
% Термины пишутся без жирного шрифта
Backbone - основная часть модели, выполняющая главное преобразование последовательности токенов\\[0.5em]
k-мер - подпоследовательность ДНК фиксированной длины $k$\\[0.5em]
Мотив - короткий консервативный паттерн в ДНК, связанный с функциональным элементом\\[0.5em]
Нуклеотид - мономерная единица ДНК (аденин A, тимин T, цитозин C, гуанин G)\\[0.5em]
Обратный комплемент - последовательность, полученная заменой каждого нуклеотида комплементарным и разворотом в обратном направлении\\[0.5em]
Спан - непрерывный диапазон позиций в исходной последовательности, покрытый одним токеном после слияния\\[0.5em]
Токенизация - процесс разбиения входной последовательности на дискретные единицы (токены) для обработки моделью\\[0.5em]
Токенизатор - модуль, преобразующий входную последовательность в последовательность токенов\\[0.5em]
Функция потерь - целевая функция, минимизируемая в процессе обучения модели
